\scanpage{338}
\mspage{38}

plus souvent purement mécaniques, mais par moments aussi créatrices) de la psyché à ceux-ci. Aussi les relations de l'enfant en bas âge à ses parents, dès avant même que les conditionnements propres à son présent milieu familial aient pu s'imprimer dans la psyché, sont dominées par ces figures archétypes ``factices'' de la Mère, de l'Epouse, du Père, et de l'Epoux qu'il a apportées avec lui ; tels des bagages lourds et encombrants, qu'il lui appartiendra d'alléger au cours de sa présente existence. Dès les premières années de l'enfance, ces archétypes factices se modifient dans un sens ou un autre sous l'influence du milieu parental et surtout, sous celle de la relation des deux parents à l'enfant. C'est dans ces années-là qu'ils prennent leur visage (le plus souvent) définitif, lequel reste plus ou moins immuable jusque dans la vieillesse. Durant tout l'Age de l'Errance, ce n'est que très exceptionnellement, dans une vie d'une fertilité spirituelle hors du commun, que l'un ou l'autre de ces visages s'affine et se décante, pour s'approcher de la Réalité.

En raison de l'importance cruciale des archétypes de la Mère et du Père, et de ceux de l'Epouse et de l'Epoux dans la vie de la psyché, \textit{notre relation à chacun de nos deux parents, et plus encore} (presque toujours) \textit{ce que fut cette relation pendant l'enfance, domine nos actes et nos comportements au cours de toute notre existence.} Il est bien entendu qu'il s'agit ici de la relation aux parents telle qu'elle est réellement, c'est-à-dire surtout de ce qui se passe dans l'inconscient moyen, et nullement de l'image toute plate et conventionnelle que nous nous plaisons à entretenir de cette relation.

De là découle d'une part, \textit{à quel point le rôle de parent est essentiel, à quel point la façon dont il est assumé ou négligé marque la vie entière de celui qui comme enfant fut confié entre nos mains.} Aussi parmi nos relations à autrui, c'est celle de parent à enfant qui, aux yeux de Dieu, implique la responsabilité la plus essentielle. C'est une responsabilité qui n'est pas éteinte par la mort de celui qui fut notre enfant. Elle accompagne celui-ci de naissance en naissance, jusqu'au moment de sa libération. (Par quoi s'éteint toute dette qu'une âme pouvait avoir à son égard...)

Mais d'autre part (et dans l'optique ``opposée'', de celui qui fut enfant), il découle de ceci à quel point il est essentiel pour l'âme, pour accéder à une maîtrise réelle de ses actes et à une véritable authenticité, d'arriver à résoudre son ``problème aux parents''. Cette tâche est inséparable de celle de clarifier notre relation tant à la mère qu'au père, en nous libérant des dépendances inconscientes par rapport à elle et à lui. De plus, en raison de la relation très forte qui relie l'un et l'autre parent non seulement aux archétypes parentaux, mais aussi à ceux du Conjoint, être libéré\deleted{e} (par Acte de Dieu) de sa dépendance par rapport aux Parents implique aussi, presque toujours, un affaiblissement de la dépendance aux archétypes de l'Epouse et de l'Epoux. Aussi une fois libérée de la ``guerre des générations'' dans sa psyché, l'âme est désormais en meilleure posture pour travailler à mettre fin aussi à la ``guerre des sexes'' qui se joue en elle.

15. La pulsion homosexuelle est une pulsion éternelle dans tout être humain

On ne peut trop insister sur cette loi fondamentale et universelle dans la vie de la psyché, que la \textit{mère est investie d'office non seulement de l'archétype de la Mère, mais aussi de celui de l'Epouse ; et de même le père est investi aussi, à côté de l'archétype ``anodin'' du Père, de celui de l'Epoux.} Cette loi, m'est-il assuré, est toujours obéie, elle ne souffre

% archive sheet number: 338
